\section{Discussion}\label{sec:discussion}
% src: S-PAPER-5, S-PAPER-8, S-PAPER-10
The answer to Q1 is no.
At this $n$, the price is not a calibrated probability as a whole.
The observed outcome-1 rate sits below the price in \GapNegBinsMain{} of \GapBinsTotal{} bins.
The effect is small and runs one way.
The net gap is \GapNetDiffMain.
The largest per-bin gap is \LargestGapAbs{} below the price, in bin \LargestGapBin{} (\Cref{tab:bins}).
We set no size line before the results.
With the minimum-trade filter, the longshot gap points the same way as the horse-racing bias.
Without the filter, the cheapest bin does not.
\todo{owner: interpretation of why prices run high, and what the gaps mean in practice for a reader of the price}

The answer to Q2 is descriptive.
Yes, by USDC volume, the top wallets do most of the trading.
This holds in the kept markets, in both roles, far beyond the 80/20 reference.
Our top \TopOneLabel{} maker share is \MakerTopOneShare, against \PaperTopShare{} in the dataset paper.
The populations differ.
We measure concentration of volume.
We do not test who sets the prices.
The k-means clusters are a side view and carry no claim.
The wallet volume includes trades after $T-\DaysShort$ day.

The answer to Q3 is descriptive too.
Prices cluster near the ends, and trade sizes are strongly right skewed.
\todo{owner: wider interpretation or policy implications, if the course report needs them}
